% Inactive fragment: bounded global shifted Hessian mass and comparison.
% Chapter regime: compact X, smooth Hermitian omega, 2 <= m < n.
\section{Hessian mass and global comparison}
\label{draft627:mass-comparison}

The maximum inequality compares measures on complementary Borel sets.
Integrating it leaves a difference of total masses. The recursive
currents express this difference through the complex Hessians of the
background powers, extending the mass calculation of
\cref{chap:mass} to bounded admissible potentials.

For bounded admissible $t_1,\ldots,t_m$, put
\[
 \mathcal M_m(t_1,\ldots,t_m)=\int_XQ_m(t_1,\ldots,t_m),
 \qquad H_m(t)=Q_m(t,\ldots,t).
\]
Here $H_m(t)$ is a finite positive Radon measure and
$\mathcal M_m$ is the total mass of a mixed positive measure, by
\cref{hessian:bounded-shifted-positive}. For a nonempty subset
$I=\{i_1<\cdots<i_q\}\subset\{1,\ldots,m\}$, put
$i(I)=i_q$ and $I^-=I\setminus\{i_q\}$. Use $D_{I^-}$ for the
recursive current with those ordered inputs, with $D_\varnothing=1$.

\begin{lemma}\label{draft627:bounded-mass-expansion}
For every bounded admissible list,
\begin{equation}\label{draft627:eq-bounded-mass-expansion}
 \mathcal M_m(t_1,\ldots,t_m)
 =\int_X\omega^n+
 \sum_{\varnothing\ne I\subset\{1,\ldots,m\}}
 \int_X t_{i(I)}D_{I^-}\wedge\ddc(\omega^{n-|I|}).
\end{equation}
Every integral on the right is an integral of bounded canonical
values against signed measure coefficients.
\end{lemma}
\begin{proof}
Expand $Q_m$ using \eqref{eq:hessian-bounded-shifted-current}.
The term for $I$ pairs $D_I$ with the smooth form
$\omega^{n-|I|}$. By \eqref{eq:hessian-bounded-recursion},
\[
 \langle D_I,\omega^{n-|I|}\rangle
 =\langle t_{i(I)}D_{I^-},\ddc(\omega^{n-|I|})\rangle.
\]
Since $|I^-|\leq m-1$, \cref{hessian:bounded-recursive-currents}
gives measure coefficients for $D_{I^-}$; their bounded product
with $t_{i(I)}$ is well defined. Thus the displayed pairing is
exactly the integral in the sum, and is finite on compact $X$.
The empty subset contributes $\int_X\omega^n$.
\end{proof}

Fix bounded admissible $z_1,\ldots,z_{m-r}$ and $1\leq r\leq m$,
and write
\[
 \mu_r(t)=Q_m(\underbrace{t,\ldots,t}_{r},z_1,\ldots,z_{m-r}).
\]
The total-mass difference below therefore has the explicit expression
\eqref{draft627:eq-bounded-mass-expansion} for the two indicated lists.

\begin{proposition}\label{draft627:comparison-mass-defect}
If $u,v$ are bounded admissible and $w=\max\{u,v\}$, then
\[
 \int_{\{u<v\}}\mathrm d\mu_r(v)
 \leq\int_{\{u<v\}}\mathrm d\mu_r(u)
                +\int_X\mathrm d\mu_r(w)-\int_X\mathrm d\mu_r(u).
\]
\end{proposition}
\begin{proof}
Set $E=\{u<v\}$. The contact-set version of
\cref{hessian:bounded-maximum-measure} gives
\[
 \mu_r(w)\geq\mathbf1_{X\setminus E}\mu_r(u)
                     +\mathbf1_E\mu_r(v).
\]
All these measures are finite and positive. Integrating on $X$ and
using $\mu_r(u)(X\setminus E)=\mu_r(u)(X)-\mu_r(u)(E)$ gives
exactly the assertion. No boundary or closed-background hypothesis
enters this calculation.
\end{proof}

The unmodified global comparison principle holds precisely when the
powers occurring in this mass formula are $\ddc$-closed. Only the
powers $n-m,\ldots,n-1$ occur at Hessian degree $m$.

\begin{theorem}\label{draft627:global-comparison-characterization}
The following conditions are equivalent:
\begin{enumerate}
\item $\ddc(\omega^{n-q})=0$ for every $1\leq q\leq m$;
\item $\int_XH_m(u)$ is independent of the smooth admissible
      potential $u$;
\item $\mathcal M_m(t_1,\ldots,t_m)=\int_X\omega^n$ for every
      bounded admissible list;
\item every pair of bounded admissible $u,v$ satisfies
\[
 \int_{\{u<v\}}H_m(v)\leq\int_{\{u<v\}}H_m(u).
\]
\end{enumerate}
Under these conditions, for every $r$ and every fixed bounded
admissible list $z_1,\ldots,z_{m-r}$, one also has
\[
 \int_{\{u<v\}}\mathrm d\mu_r(v)
           \leq\int_{\{u<v\}}\mathrm d\mu_r(u).
\]
\end{theorem}
\begin{proof}
The first condition implies the third by
\eqref{draft627:eq-bounded-mass-expansion}. The third implies the
second by taking all inputs equal to a smooth potential.

To prove that the second implies the first, fix any smooth real
$f$. For all sufficiently small real $t$ of either sign,
$\omega+t\ddc f$ is positive, so $tf$ is strictly admissible.
The polynomial
\[
 \int_X(\omega+t\ddc f)^m\wedge\omega^{n-m}
 =\sum_{q=0}^m\binom mq t^q
              \int_X(\ddc f)^q\wedge\omega^{n-q}
\]
is constant on an interval. Consequently every coefficient with
$q\geq1$ vanishes. Polarizing that coefficient by replacing $f$
with $\sum_{a=1}^q s_af_a$ gives
\[
 \int_X\ddc f_1\wedge\cdots\wedge\ddc f_q
                        \wedge\omega^{n-q}=0
\]
for all smooth real $f_a$. Integration by parts in the first
factor, using \cref{lem:ibp}, gives
\[
 \int_Xf_1\ddc f_2\wedge\cdots\wedge\ddc f_q
                         \wedge\ddc(\omega^{n-q})=0.
\]
Since $f_1$ is arbitrary, the smooth top-degree form paired with
it vanishes pointwise. At each point the other complex Hessians
can be prescribed by real quadratic functions times a cutoff.
Complexifying this multilinear identity gives all coordinate
$(q-1,q-1)$-forms, as in the proof of
\cref{thm:mass-invariance-characterization}. They detect every
coefficient of $\ddc(\omega^{n-q})$, so that form vanishes. When
$q=1$, there are no other Hessians and the same conclusion follows
directly from the arbitrariness of $f_1$.

Under the third condition, $\mu_r(w)$ and $\mu_r(u)$ have the same
total mass. \Cref{draft627:comparison-mass-defect} proves the mixed
comparison assertion, and $r=m$ gives the fourth condition.
Finally, suppose the fourth condition holds and fix smooth admissible
$u,v$. Choose $C$ so large that $u<v+C$ everywhere.
Constant invariance gives $H_m(v+C)=H_m(v)$, so comparison gives
$\int_XH_m(v)\leq\int_XH_m(u)$. Exchanging $u,v$ gives equality,
which is the second condition.
\end{proof}

The same varying background that obstructs Monge--Amp\`ere comparison
also obstructs comparison at each intermediate Hessian degree.

\begin{example}\label{draft627:torus-comparison-obstruction}
Let $X=(\mathbb C/(\mathbb Z+\mathrm i\mathbb Z))^n$, with
$z_j=x_j+\mathrm i y_j$, and put
\[
 \beta_j=\frac{\mathrm i}{2\pi}\,\mathrm dz_j\wedge\mathrm d\bar z_j,
 \qquad b=2+\cos(2\pi x_1),
 \qquad\omega=\beta_1+b\beta_2+\sum_{j=3}^n\beta_j.
\]
For $|t|<\pi^{-2}$ the potential $u_t=t\cos(2\pi x_1)$ is smooth
strictly admissible, and, with $B_n=\int_X\beta_1\wedge\cdots\wedge\beta_n$,
\[
 \int_X H_m(u_t)=n!\left(2-\frac{m\pi^2t}{2n}\right)B_n.
\]
Thus every $-\pi^{-2}<t<0$, together with a constant $C>|t|$,
gives a pair $u=0$, $v=u_t+C$ that violates the unmodified global
comparison inequality.
\end{example}
\begin{proof}
As in \cref{ex:torus-mass},
$\ddc u_t=-\pi^2t\cos(2\pi x_1)\beta_1$.
The coefficient of $\beta_1$ in $\omega+\ddc u_t$ is strictly
positive in the stated range, and all the other coefficients are
positive. Since $(\ddc u_t)^2=0$, expansion gives
\[
 H_m(u_t)=\omega^n+m\ddc u_t\wedge\omega^{n-1}
 =n!b\left(1-\frac mn\pi^2t\cos(2\pi x_1)\right)
                         \beta_1\wedge\cdots\wedge\beta_n.
\]
The period averages of $b$ and $b\cos(2\pi x_1)$ are respectively
$2$ and $1/2$, proving the mass formula. For $t<0$ this mass is
strictly larger than $\int_XH_m(0)=2n!B_n$. Meanwhile $C>|t|$
gives $\{0<u_t+C\}=X$, so the comparison inequality for this pair
would assert the opposite ordering of these masses.
\end{proof}
